% Source: https://artofproblemsolving.com/community/c3h3259043_video_solutions_to_2024_math_in_the_matrix_competition_blitz_round__discussion
% Stable identity: aops-topic-3259043-q07
After the first step each is at one. After the second each is at zero or two with probabilities $\dfrac{1}{2}$. Therefore after the third each is at one with probability $\dfrac{1}{2}+\dfrac{1}{4}=\dfrac{3}{4}$, and at three with probability $\dfrac{1}{4}$. Independence gives coincidence probability $(\dfrac{3}{4})^2+(\dfrac{1}{4})^2=\dfrac{5}{8}$. Thus $\boxed{m+n=13}$.
